
\begin{table}[H]\centering\small
\caption{Levelised cost per kL of design yield (2024 rupees) by group, and sensitivity.}\label{tab:coststats}
\begin{tabular}{lrrr}\toprule
Group & n & Median Rs/kL & Range\\ \midrule
All & 27 & 20.9 (95\% CI 8.5--43.3) & IQR 6.9--65.9\\
individual house & 3 & 39.8 & 28.8--43.3\\
residential colony & 6 & 15.4 & 1.8--98.3\\
institutional & 15 & 15.2 & 1.4--198.5\\
industrial & 3 & 110.7 & 6.8--151.7\\
system: recharge & 22 & 14.6 & \\
system: storage & 3 & 43.3 & \\
system: both & 2 & 118.7 & \\
CSE Delhi-NCR rows (2000--2004) & 20 & 13.4 &\\
Rows with cost year $\geq$ 2010 & 5 & 98.3 &\\
Sensitivity: 20-year life, 6\% discount & 27 & 16.4 &\\
Sensitivity: 15\% discount (the CGWB example rate) & 27 & 29.2 &\\
Rows with cost year $\geq$ 2010 at observed yield & 5 & 199 &\\
At observed yield (multiplier 2.02) & 27 & 42 &\\
\bottomrule\end{tabular}
\end{table}
Kruskal-Wallis across the four building-scale groups: H = 2.07, p = 0.559. Spearman $\rho$ of $\log_{10}$ area against Rs/kL: -0.74 (p = 0.0000, n = 26); against capital per m$^2$: -0.73 (p = 0.0000). The delivered-yield multiplier is the design-calculated median ratio (design yield as a share of potential, 0.56) divided by the observed median (0.28).
